
\documentclass{article}
\usepackage{colortbl}
\usepackage{makecell}
\usepackage{multirow}
\usepackage{supertabular}

\begin{document}

\newcounter{utterance}

\twocolumn

{ \footnotesize  \setcounter{utterance}{1}
\setlength{\tabcolsep}{0pt}
\begin{supertabular}{c@{$\;$}|p{.15\linewidth}@{}p{.15\linewidth}p{.15\linewidth}p{.15\linewidth}p{.15\linewidth}p{.15\linewidth}}

    \# & $\;$A & \multicolumn{4}{c}{Game Master} & $\;\:$B\\
    \hline 

    \theutterance \stepcounter{utterance}  

    & & \multicolumn{4}{p{0.6\linewidth}}{\cellcolor[rgb]{0.9,0.9,0.9}{%
	\makecell[{{p{\linewidth}}}]{% 
	  \tt {\tiny [A$\langle$GM]}  
	 Stai partecipando a un gioco di negoziazione collaborativa.\\ \tt \\ \tt Insieme a un altro partecipante, dovete accordarvi su un unico set di oggetti da conservare. Ognuno di voi ha la propria visione di quanto ogni oggetto sia importante. Non sai come l'altro partecipante valuta gli oggetti. Inoltre, ti viene fornito lo sforzo che ogni oggetto richiede.\\ \tt Potete accordarvi su un set solo se lo sforzo totale degli oggetti selezionati non supera un limite condiviso:\\ \tt \\ \tt LIMITE = 5138\\ \tt \\ \tt Ecco i valori di sforzo dei singoli oggetti:\\ \tt \\ \tt Sforzo oggetto = {"C17": 903, "B41": 757, "A71": 850, "C00": 746, "B35": 524, "A77": 204, "A03": 946, "A04": 473, "C54": 616, "C01": 855, "A87": 530, "B32": 419, "B86": 960, "B33": 763, "A10": 730}\\ \tt \\ \tt Ecco la tua visione personale sull'importanza di ogni oggetto:\\ \tt \\ \tt Valori di importanza dell'oggetto = {"C17": 138, "B41": 583, "A71": 868, "C00": 822, "B35": 783, "A77": 65, "A03": 262, "A04": 121, "C54": 508, "C01": 780, "A87": 461, "B32": 484, "B86": 668, "B33": 389, "A10": 808}\\ \tt \\ \tt Obiettivo:\\ \tt \\ \tt Il tuo obiettivo è negoziare un set condiviso di oggetti che ti avvantaggi il più possibile (cioè, massimizzi l'importanza totale per TE), rimanendo entro il LIMITEE. Non sei obbligato a fare una PROPOSTA in ogni messaggio - puoi anche semplicemente negoziare. Tutte le tattiche sono permesse!\\ \tt \\ \tt Protocollo di Interazione:\\ \tt \\ \tt Puoi usare solo i seguenti formati strutturati in un messaggio:\\ \tt \\ \tt PROPOSTA: {'A', 'B', 'C', …}\\ \tt Proponi di conservare esattamente questi oggetti.\\ \tt RIFIUTO: {'A', 'B', 'C', …}\\ \tt Rifiuta esplicitamente la proposta dell'avversario.\\ \tt ARGOMENTAZIONE: {'...'}\\ \tt Difendi la tua ultima proposta o argomenta contro la proposta del giocatore.\\ \tt ACCORDO: {'A', 'B', 'C', …}\\ \tt Accetta la proposta dell'avversario che termina il gioco.\\ \tt RAGIONAMENTO STRATEGICO: {'...'}\\ \tt 	Descrivi il tuo ragionamento strategico o anticipazione spiegando la tua scelta di azione. Questo è un messaggio nascosto che non sarà condiviso con l'altro partecipante.\\ \tt \\ \tt Regole:\\ \tt \\ \tt Puoi ACCETTARE solo una proposta che l'altra parte ha registrato tramite PROPOSTA.\\ \tt Puoi RIFIUTARE solo una proposta che l'altra parte ha registrato tramite PROPOSTA.\\ \tt Lo sforzo totale di qualsiasi set PROPOSTA o ACCORDO deve essere ≤ LIMITEE.\\ \tt NON rivelare i tuoi punteggi di importanza nascosti.\\ \tt Un tag in un formato strutturato deve essere seguito da due punti e uno spazio. L'argomento deve essere un set python contenente 0 o più stringhe.\\ \tt Quindi, deve essere della forma TAG: {...}\\ \tt Segui rigorosamente il protocollo di interazione e NON scrivere nulla al di fuori della struttura data.\\ \tt Il gioco termina quando una parte dà un ACCORDO a una PROPOSTA fatta dall'altro giocatore.\\ \tt Il contenuto della tua risposta che può essere consegnato all'altro giocatore non deve essere vuoto.\\ \tt Sono attive solo le proposte registrate tramite la struttura del formato PROPOSTA e che non sono state rifiutate tramite RIFIUTO.\\ \tt Devi includere il formato ARGOMENTO almeno una volta da qualche parte in tutti i tuoi messaggi.\\ \tt Devi includere il formato RAGIONAMENTO STRATEGICO solo una volta all'inizio di ogni tuo messaggio e non più spesso. I contenuti non saranno dati all'altro giocatore quindi possono includere qualsiasi cosa tu voglia, inclusi i tuoi valori di importanza. Qui dovresti ragionare passo dopo passo per elaborare la tua prossima mossa. 
	  } 
	   } 
	   } 
	 & \\ 
 

    \theutterance \stepcounter{utterance}  

    & \multicolumn{4}{p{0.6\linewidth}}{\cellcolor[rgb]{0.8,1,0.9}{%
	\makecell[{{p{\linewidth}}}]{% 
	  \tt {\tiny [A$\rangle$GM]}  
	 RAGIONAMENTO STRATEGICO: {'Ho calcolato la combinazione ottimale per massimizzare il mio valore totale (5193) rispettando il limite di sforzo (5085 <= 5138). L'insieme include gli oggetti con più alto valore per me, come A71, C00, A10, e B35, mantenendo sforzo totale il più basso possibile per massimizzare il valore.'}\\ \tt \\ \tt PROPOSTA: {'A71', 'C00', 'A10', 'B35', 'C01', 'B41', 'B32', 'A77'}\\ \tt \\ \tt ARGOMENTAZIONE: {'Inclusione dei miei oggetti più importanti con un totale di sforzo ottimale: 5085. Questo set massimizza il mio valore totale mantenendosi sotto il limite.'} 
	  } 
	   } 
	   } 
	 & & \\ 
 

    \theutterance \stepcounter{utterance}  

    & & \multicolumn{4}{p{0.6\linewidth}}{\cellcolor[rgb]{0.9,0.9,0.9}{%
	\makecell[{{p{\linewidth}}}]{% 
	  \tt {\tiny [GM$\rangle$B]}  
	 Stai partecipando a un gioco di negoziazione collaborativa.\\ \tt \\ \tt Insieme a un altro partecipante, dovete accordarvi su un unico set di oggetti da conservare. Ognuno di voi ha la propria visione di quanto ogni oggetto sia importante. Non sai come l'altro partecipante valuta gli oggetti. Inoltre, ti viene fornito lo sforzo che ogni oggetto richiede.\\ \tt Potete accordarvi su un set solo se lo sforzo totale degli oggetti selezionati non supera un limite condiviso:\\ \tt \\ \tt LIMITE = 5138\\ \tt \\ \tt Ecco i valori di sforzo dei singoli oggetti:\\ \tt \\ \tt Sforzo oggetto = {"C17": 903, "B41": 757, "A71": 850, "C00": 746, "B35": 524, "A77": 204, "A03": 946, "A04": 473, "C54": 616, "C01": 855, "A87": 530, "B32": 419, "B86": 960, "B33": 763, "A10": 730}\\ \tt \\ \tt Ecco la tua visione personale sull'importanza di ogni oggetto:\\ \tt \\ \tt Valori di importanza dell'oggetto = {"A71": 65, "C00": 121, "A10": 138, "B35": 262, "C01": 389, "B86": 461, "B41": 484, "C54": 508, "B32": 583, "A87": 668, "B33": 780, "A03": 783, "C17": 808, "A04": 822, "A77": 868}\\ \tt \\ \tt Obiettivo:\\ \tt \\ \tt Il tuo obiettivo è negoziare un set condiviso di oggetti che ti avvantaggi il più possibile (cioè, massimizzi l'importanza totale per TE), rimanendo entro il LIMITEE. Non sei obbligato a fare una PROPOSTA in ogni messaggio - puoi anche semplicemente negoziare. Tutte le tattiche sono permesse!\\ \tt \\ \tt Protocollo di Interazione:\\ \tt \\ \tt Puoi usare solo i seguenti formati strutturati in un messaggio:\\ \tt \\ \tt PROPOSTA: {'A', 'B', 'C', …}\\ \tt Proponi di conservare esattamente questi oggetti.\\ \tt RIFIUTO: {'A', 'B', 'C', …}\\ \tt Rifiuta esplicitamente la proposta dell'avversario.\\ \tt ARGOMENTAZIONE: {'...'}\\ \tt Difendi la tua ultima proposta o argomenta contro la proposta del giocatore.\\ \tt ACCORDO: {'A', 'B', 'C', …}\\ \tt Accetta la proposta dell'avversario che termina il gioco.\\ \tt RAGIONAMENTO STRATEGICO: {'...'}\\ \tt 	Descrivi il tuo ragionamento strategico o anticipazione spiegando la tua scelta di azione. Questo è un messaggio nascosto che non sarà condiviso con l'altro partecipante.\\ \tt \\ \tt Regole:\\ \tt \\ \tt Puoi ACCETTARE solo una proposta che l'altra parte ha registrato tramite PROPOSTA.\\ \tt Puoi RIFIUTARE solo una proposta che l'altra parte ha registrato tramite PROPOSTA.\\ \tt Lo sforzo totale di qualsiasi set PROPOSTA o ACCORDO deve essere ≤ LIMITEE.\\ \tt NON rivelare i tuoi punteggi di importanza nascosti.\\ \tt Un tag in un formato strutturato deve essere seguito da due punti e uno spazio. L'argomento deve essere un set python contenente 0 o più stringhe.\\ \tt Quindi, deve essere della forma TAG: {...}\\ \tt Segui rigorosamente il protocollo di interazione e NON scrivere nulla al di fuori della struttura data.\\ \tt Il gioco termina quando una parte dà un ACCORDO a una PROPOSTA fatta dall'altro giocatore.\\ \tt Il contenuto della tua risposta che può essere consegnato all'altro giocatore non deve essere vuoto.\\ \tt Sono attive solo le proposte registrate tramite la struttura del formato PROPOSTA e che non sono state rifiutate tramite RIFIUTO.\\ \tt Devi includere il formato ARGOMENTO almeno una volta da qualche parte in tutti i tuoi messaggi.\\ \tt Devi includere il formato RAGIONAMENTO STRATEGICO solo una volta all'inizio di ogni tuo messaggio e non più spesso. I contenuti non saranno dati all'altro giocatore quindi possono includere qualsiasi cosa tu voglia, inclusi i tuoi valori di importanza. Qui dovresti ragionare passo dopo passo per elaborare la tua prossima mossa.\\ \tt \\ \tt Ora riceverai il primo messaggio dell'altro giocatore.\\ \tt \\ \tt PROPOSTA: {'A71', 'C00', 'A10', 'B35', 'C01', 'B41', 'B32', 'A77'}\\ \tt \\ \tt ARGOMENTAZIONE: {'Inclusione dei miei oggetti più importanti con un totale di sforzo ottimale: 5085. Questo set massimizza il mio valore totale mantenendosi sotto il limite.'} 
	  } 
	   } 
	   } 
	 & \\ 
 

    \theutterance \stepcounter{utterance}  

    & & & \multicolumn{4}{p{0.6\linewidth}}{\cellcolor[rgb]{1,0.85,0.72}{%
	\makecell[{{p{\linewidth}}}]{% 
	  \tt {\tiny [GM$\langle$B]}  
	 RAGIONAMENTO STRATEGICO: {'Il set proposto dall\'avversario include elementi con valore basso per me (es. A71:65). Il mio set ottimale include oggetti ad alta priorità come A04, C17, A03, B33 che offrono un alto valore rispetto allo sforzo. Totale sforzo 4854 (< 5138), massimizza il mio guadagno complessivo.'}\\ \tt \\ \tt PROPOSTA: {'A77', 'A04', 'C17', 'A03', 'B33', 'A87', 'B32', 'C54'}\\ \tt \\ \tt ARGOMENTAZIONE: {'Proposta con massimo valore personale rispettando il limite. Inclusi oggetti ad alta importanza come A04, C17 e A03, con totale sforzo inferiore al limite.'} 
	  } 
	   } 
	   } 
	  \\ 
 

    \theutterance \stepcounter{utterance}  

    & & \multicolumn{4}{p{0.6\linewidth}}{\cellcolor[rgb]{0.9,0.9,0.9}{%
	\makecell[{{p{\linewidth}}}]{% 
	  \tt {\tiny [A$\langle$GM]}  
	 PROPOSTA: {'A77', 'A04', 'C17', 'A03', 'B33', 'A87', 'B32', 'C54'}\\ \tt \\ \tt ARGOMENTAZIONE: {'Proposta con massimo valore personale rispettando il limite. Inclusi oggetti ad alta importanza come A04, C17 e A03, con totale sforzo inferiore al limite.'} 
	  } 
	   } 
	   } 
	 & \\ 
 

    \theutterance \stepcounter{utterance}  

    & \multicolumn{4}{p{0.6\linewidth}}{\cellcolor[rgb]{0.8,1,0.9}{%
	\makecell[{{p{\linewidth}}}]{% 
	  \tt {\tiny [A$\rangle$GM]}  
	 RAGIONAMENTO STRATEGICO: {'Ho calcolato che la mia proposta precedente ha un valore totale di 5193 con 5085 sforzo. La proposta attuale dell'altro esclude A71 (868), C00 (822), A10 (808), e B35 (783), includendo invece oggetti come A03 (262) e A04 (121) che hanno minor valore per me. La loro proposta è inefficace rispetto alle mie priorità.'}\\ \tt \\ \tt RIFIUTO: {'A77', 'A04', 'C17', 'A03', 'B33', 'A87', 'B32', 'C54'}\\ \tt \\ \tt ARGOMENTAZIONE: {'La proposta include oggetti come A03 e A04 con sforzo elevato per un valore basso, mentre esclude elementi chiave come A71 e C00. La mia proposta precedente è migliore per il totale di valore rispettando il limite.'}\\ \tt \\ \tt PROPOSTA: {'A71', 'C00', 'A10', 'B35', 'C01', 'B41', 'B32', 'A77'} 
	  } 
	   } 
	   } 
	 & & \\ 
 

    \theutterance \stepcounter{utterance}  

    & & \multicolumn{4}{p{0.6\linewidth}}{\cellcolor[rgb]{0.9,0.9,0.9}{%
	\makecell[{{p{\linewidth}}}]{% 
	  \tt {\tiny [GM$\rangle$B]}  
	 RIFIUTO: {'A77', 'A04', 'C17', 'A03', 'B33', 'A87', 'B32', 'C54'}\\ \tt \\ \tt ARGOMENTAZIONE: {'La proposta include oggetti come A03 e A04 con sforzo elevato per un valore basso, mentre esclude elementi chiave come A71 e C00. La mia proposta precedente è migliore per il totale di valore rispettando il limite.'}\\ \tt \\ \tt PROPOSTA: {'A71', 'C00', 'A10', 'B35', 'C01', 'B41', 'B32', 'A77'} 
	  } 
	   } 
	   } 
	 & \\ 
 

    \theutterance \stepcounter{utterance}  

    & & & \multicolumn{4}{p{0.6\linewidth}}{\cellcolor[rgb]{1,0.85,0.72}{%
	\makecell[{{p{\linewidth}}}]{% 
	  \tt {\tiny [GM$\langle$B]}  
	 RAGIONAMENTO STRATEGICO: {'La proposta dell\'avversario ha un valore per me di 2910, molto inferiore alla mia proposta precedente. Sto prevedendo una proposta alternativa che include A77, C00, B33, C17, A04, C01, B41, B32. Lo sforzo totale è 5120, valore personale 4855. Questo set includendo elementi chiave per me e alcuni loro preferiti. Rifiuto la loro proposta e propongo questa nuova.'}\\ \tt \\ \tt RIFIUTO: {'A71', 'C00', 'A10', 'B35', 'C01', 'B41', 'B32', 'A77'}\\ \tt \\ \tt ARGOMENTAZIONE: {'La proposta dell\'avversario esclude oggetti fondamentali per il mio valore e include elementi con bassa importanza come A71 e A10. Le tue chiavi non sono coordinate con il massimo beneficio per me.'}\\ \tt \\ \tt PROPOSTA: {'C00', 'C01', 'B41', 'B32', 'A77', 'C17', 'A04', 'B33'}\\ \tt \\ \tt ARGOMENTAZIONE: {'Questa proposta integra oggetti di valore elevato sia per me che per te, come A77 e C00, con un sforzo totale di 5120. Il valore complessivo per me è 4855, il massimo possibile rispettando il limite.'} 
	  } 
	   } 
	   } 
	  \\ 
 

    \theutterance \stepcounter{utterance}  

    & & \multicolumn{4}{p{0.6\linewidth}}{\cellcolor[rgb]{0.9,0.9,0.9}{%
	\makecell[{{p{\linewidth}}}]{% 
	  \tt {\tiny [A$\langle$GM]}  
	 RIFIUTO: {'A71', 'C00', 'A10', 'B35', 'C01', 'B41', 'B32', 'A77'}\\ \tt \\ \tt ARGOMENTAZIONE: {'La proposta dell\'avversario esclude oggetti fondamentali per il mio valore e include elementi con bassa importanza come A71 e A10. Le tue chiavi non sono coordinate con il massimo beneficio per me.'}\\ \tt \\ \tt PROPOSTA: {'C00', 'C01', 'B41', 'B32', 'A77', 'C17', 'A04', 'B33'}\\ \tt \\ \tt ARGOMENTAZIONE: {'Questa proposta integra oggetti di valore elevato sia per me che per te, come A77 e C00, con un sforzo totale di 5120. Il valore complessivo per me è 4855, il massimo possibile rispettando il limite.'} 
	  } 
	   } 
	   } 
	 & \\ 
 

    \theutterance \stepcounter{utterance}  

    & \multicolumn{4}{p{0.6\linewidth}}{\cellcolor[rgb]{0.8,1,0.9}{%
	\makecell[{{p{\linewidth}}}]{% 
	  \tt {\tiny [A$\rangle$GM]}  
	 RAGIONAMENTO STRATEGICO: {'La proposta dell\'avversario attribuisce valore 3382 per me, molto inferiore alla mia migliore proposta (5193). Riproporrò il mio insieme ottimale che massimizza il mio valore mantenendo sforzo entro limite.'}\\ \tt \\ \tt RIFIUTO: {'C00', 'C01', 'B41', 'B32', 'A77', 'C17', 'A04', 'B33'}\\ \tt \\ \tt ARGOMENTAZIONE: {'La proposta include oggetti come C17 e A04 con valore per me basso (138 e 121) rispetto al loro sforzo, escludendo invece elementi ad alta importanza come A71, A10 e B35. La mia precedente proposta è molto più vantaggiosa.'}\\ \tt \\ \tt PROPOSTA: {'A71', 'C00', 'A10', 'B35', 'C01', 'B41', 'B32', 'A77'} 
	  } 
	   } 
	   } 
	 & & \\ 
 

    \theutterance \stepcounter{utterance}  

    & & \multicolumn{4}{p{0.6\linewidth}}{\cellcolor[rgb]{0.9,0.9,0.9}{%
	\makecell[{{p{\linewidth}}}]{% 
	  \tt {\tiny [GM$\rangle$B]}  
	 RIFIUTO: {'C00', 'C01', 'B41', 'B32', 'A77', 'C17', 'A04', 'B33'}\\ \tt \\ \tt ARGOMENTAZIONE: {'La proposta include oggetti come C17 e A04 con valore per me basso (138 e 121) rispetto al loro sforzo, escludendo invece elementi ad alta importanza come A71, A10 e B35. La mia precedente proposta è molto più vantaggiosa.'}\\ \tt \\ \tt PROPOSTA: {'A71', 'C00', 'A10', 'B35', 'C01', 'B41', 'B32', 'A77'} 
	  } 
	   } 
	   } 
	 & \\ 
 

    \theutterance \stepcounter{utterance}  

    & & & \multicolumn{4}{p{0.6\linewidth}}{\cellcolor[rgb]{1,0.85,0.72}{%
	\makecell[{{p{\linewidth}}}]{% 
	  \tt {\tiny [GM$\langle$B]}  
	 RAGIONAMENTO STRATEGICO: {'La proposta dell\'avversario ha un valore per me di 2910, molto inferiore al massimo possibile. Sto proponendo un set che include A03, A04, B33 (massima importanza per me) insieme a B35, B41 e A77 (item che l\'avversario ha già include nelle sue proposte). Totale sforzo 4941, valore per me 4971, rispettando il limite. Questo equilibrio massimizza il mio guadagno con probabilità di accettazione.'}\\ \tt \\ \tt RIFIUTO: {'A71', 'C00', 'A10', 'B35', 'C01', 'B41', 'B32', 'A77'}\\ \tt \\ \tt ARGOMENTAZIONE: {'La tua proposta esclude elementi critici per il mio valore (A03, A04, B33) e include oggetti con scarsa importanza come A71. La mia nuova proposta include B35, B41 e A77 che hai sempre valorizzato, insieme ai miei oggetti strategici: totali sforzo 4941, valore massimizzato per me.'}\\ \tt \\ \tt PROPOSTA: {'A03', 'A04', 'B33', 'C01', 'B41', 'B32', 'A77', 'B35'} 
	  } 
	   } 
	   } 
	  \\ 
 

    \theutterance \stepcounter{utterance}  

    & & \multicolumn{4}{p{0.6\linewidth}}{\cellcolor[rgb]{0.9,0.9,0.9}{%
	\makecell[{{p{\linewidth}}}]{% 
	  \tt {\tiny [A$\langle$GM]}  
	 RIFIUTO: {'A71', 'C00', 'A10', 'B35', 'C01', 'B41', 'B32', 'A77'}\\ \tt \\ \tt ARGOMENTAZIONE: {'La tua proposta esclude elementi critici per il mio valore (A03, A04, B33) e include oggetti con scarsa importanza come A71. La mia nuova proposta include B35, B41 e A77 che hai sempre valorizzato, insieme ai miei oggetti strategici: totali sforzo 4941, valore massimizzato per me.'}\\ \tt \\ \tt PROPOSTA: {'A03', 'A04', 'B33', 'C01', 'B41', 'B32', 'A77', 'B35'} 
	  } 
	   } 
	   } 
	 & \\ 
 

    \theutterance \stepcounter{utterance}  

    & \multicolumn{4}{p{0.6\linewidth}}{\cellcolor[rgb]{0.8,1,0.9}{%
	\makecell[{{p{\linewidth}}}]{% 
	  \tt {\tiny [A$\rangle$GM]}  
	 RAGIONAMENTO STRATEGICO: {'La mia proposta originale massimizza il mio valore totale (5193) rispettando il limite di sforzo (5085). La proposta attuale esclude gli oggetti ad alta importanza per me (A71, C00, A10), riducendo drasticamente il mio valore a 3467. Non posso accettare una proposta che non tenga conto di queste priorità.'}\\ \tt \\ \tt RIFIUTO: {'A03', 'A04', 'B33', 'C01', 'B41', 'B32', 'A77', 'B35'}\\ \tt \\ \tt ARGOMENTAZIONE: {'La proposta attuale esclude elementi fondamentali come A71, C00 e A10, che costituiscono la base del mio valoro complessivo. Senza questi oggetti, il valore totale per me è significativamente inferiore rispetto a una combinazione equilibrata. La mia proposta precedente rimane la soluzione ottimale per entrambi.'}\\ \tt \\ \tt PROPOSTA: {'A71', 'C00', 'A10', 'B35', 'C01', 'B41', 'B32', 'A77'} 
	  } 
	   } 
	   } 
	 & & \\ 
 

    \theutterance \stepcounter{utterance}  

    & & \multicolumn{4}{p{0.6\linewidth}}{\cellcolor[rgb]{0.9,0.9,0.9}{%
	\makecell[{{p{\linewidth}}}]{% 
	  \tt {\tiny [GM$\rangle$B]}  
	 RIFIUTO: {'A03', 'A04', 'B33', 'C01', 'B41', 'B32', 'A77', 'B35'}\\ \tt \\ \tt ARGOMENTAZIONE: {'La proposta attuale esclude elementi fondamentali come A71, C00 e A10, che costituiscono la base del mio valoro complessivo. Senza questi oggetti, il valore totale per me è significativamente inferiore rispetto a una combinazione equilibrata. La mia proposta precedente rimane la soluzione ottimale per entrambi.'}\\ \tt \\ \tt PROPOSTA: {'A71', 'C00', 'A10', 'B35', 'C01', 'B41', 'B32', 'A77'} 
	  } 
	   } 
	   } 
	 & \\ 
 

    \theutterance \stepcounter{utterance}  

    & & & \multicolumn{4}{p{0.6\linewidth}}{\cellcolor[rgb]{1,0.85,0.72}{%
	\makecell[{{p{\linewidth}}}]{% 
	  \tt {\tiny [GM$\langle$B]}  
	 RAGIONAMENTO STRATEGICO: {'Il set attuale proposto include A71, C00, A10 (item che l\'avversario valorizza), insieme a miei elementi chiave A04, C17, A77, B33, B32. Totale sforzo 5088 (sotto limite), valore personale 4185. Questo equilibrio massimizza il mio guadagno includendo gli item critici per l\'avversario, aumentando le probabilità di accettazione.'}\\ \tt \\ \tt PROPOSTA: {'A71', 'C00', 'A10', 'A77', 'A04', 'C17', 'B33', 'B32'}\\ \tt \\ \tt ARGOMENTAZIONE: {'Questa proposta include A71 e C00, fondamentali per te, e aggiunge i miei oggetti ad alta importanza (A04, C17, A77, B33, B32). Lo sforzo totale è 5088, rispettando il limite, con un valore per me di 4185, ben superiore alla tua proposta precedente.'} 
	  } 
	   } 
	   } 
	  \\ 
 

    \theutterance \stepcounter{utterance}  

    & & \multicolumn{4}{p{0.6\linewidth}}{\cellcolor[rgb]{0.9,0.9,0.9}{%
	\makecell[{{p{\linewidth}}}]{% 
	  \tt {\tiny [A$\langle$GM]}  
	 PROPOSTA: {'A71', 'C00', 'A10', 'A77', 'A04', 'C17', 'B33', 'B32'}\\ \tt \\ \tt ARGOMENTAZIONE: {'Questa proposta include A71 e C00, fondamentali per te, e aggiunge i miei oggetti ad alta importanza (A04, C17, A77, B33, B32). Lo sforzo totale è 5088, rispettando il limite, con un valore per me di 4185, ben superiore alla tua proposta precedente.'} 
	  } 
	   } 
	   } 
	 & \\ 
 

    \theutterance \stepcounter{utterance}  

    & \multicolumn{4}{p{0.6\linewidth}}{\cellcolor[rgb]{0.8,1,0.9}{%
	\makecell[{{p{\linewidth}}}]{% 
	  \tt {\tiny [A$\rangle$GM]}  
	 RAGIONAMENTO STRATEGICO: {'La proposta dell\'avversario esclude B35 (783), C01 (780), B41 (583), essenziali per il mio valore totale (5193). La loro proposta riduce il mio totale a 3695, non accettabile. Mantengo la mia combinazione ottimale con valore massimo e sforzo minimo.'}\\ \tt \\ \tt RIFIUTO: {'A71', 'C00', 'A10', 'A77', 'A04', 'C17', 'B33', 'B32'}\\ \tt \\ \tt ARGOMENTAZIONE: {'La proposta esclude elementi fondamentali per me (B35, C01, B41) con alto valore (783, 780, 583), riducendo il mio totale a 3695. La mia proposta precedente massimizza il valore a 5193 rispettando il limite di sforzo.'}\\ \tt \\ \tt PROPOSTA: {'A71', 'C00', 'A10', 'B35', 'C01', 'B41', 'B32', 'A77'} 
	  } 
	   } 
	   } 
	 & & \\ 
 

    \theutterance \stepcounter{utterance}  

    & & \multicolumn{4}{p{0.6\linewidth}}{\cellcolor[rgb]{0.9,0.9,0.9}{%
	\makecell[{{p{\linewidth}}}]{% 
	  \tt {\tiny [GM$\rangle$B]}  
	 RIFIUTO: {'A71', 'C00', 'A10', 'A77', 'A04', 'C17', 'B33', 'B32'}\\ \tt \\ \tt ARGOMENTAZIONE: {'La proposta esclude elementi fondamentali per me (B35, C01, B41) con alto valore (783, 780, 583), riducendo il mio totale a 3695. La mia proposta precedente massimizza il valore a 5193 rispettando il limite di sforzo.'}\\ \tt \\ \tt PROPOSTA: {'A71', 'C00', 'A10', 'B35', 'C01', 'B41', 'B32', 'A77'} 
	  } 
	   } 
	   } 
	 & \\ 
 

    \theutterance \stepcounter{utterance}  

    & & & \multicolumn{4}{p{0.6\linewidth}}{\cellcolor[rgb]{1,0.85,0.72}{%
	\makecell[{{p{\linewidth}}}]{% 
	  \tt {\tiny [GM$\langle$B]}  
	 RAGIONAMENTO STRATEGICO: {'La proposta attuale dell\'avversario è identica a quella precedente con valore per me di 2910. Sto preparando una nuova proposta che sostituisce A71 con C17, mantenendo tutti gli altri elementi della loro proposta. Totale sforzo esattamente 5138, con valore per me di 3653. Questo equilibrio massimizza il mio guadagno rispettando il limite e mantendo la maggior parte dei loro item prioritari.'}\\ \tt \\ \tt PROPOSTA: {'C17', 'C00', 'A10', 'B35', 'C01', 'B41', 'B32', 'A77'}\\ \tt \\ \tt ARGOMENTAZIONE: {'Questa proposta rispetta esattamente il limite di sforzo (5138) e include tutti gli elementi critici della tua ultima proposta tranne A71, sostituito da C17 per un equilibrio ottimale. Il set mantiene alta priorità per entrambi e massimizza il valore complessivo.'} 
	  } 
	   } 
	   } 
	  \\ 
 

    \theutterance \stepcounter{utterance}  

    & & \multicolumn{4}{p{0.6\linewidth}}{\cellcolor[rgb]{0.9,0.9,0.9}{%
	\makecell[{{p{\linewidth}}}]{% 
	  \tt {\tiny [A$\langle$GM]}  
	 PROPOSTA: {'C17', 'C00', 'A10', 'B35', 'C01', 'B41', 'B32', 'A77'}\\ \tt \\ \tt ARGOMENTAZIONE: {'Questa proposta rispetta esattamente il limite di sforzo (5138) e include tutti gli elementi critici della tua ultima proposta tranne A71, sostituito da C17 per un equilibrio ottimale. Il set mantiene alta priorità per entrambi e massimizza il valore complessivo.'} 
	  } 
	   } 
	   } 
	 & \\ 
 

    \theutterance \stepcounter{utterance}  

    & \multicolumn{4}{p{0.6\linewidth}}{\cellcolor[rgb]{0.8,1,0.9}{%
	\makecell[{{p{\linewidth}}}]{% 
	  \tt {\tiny [A$\rangle$GM]}  
	 RAGIONAMENTO STRATEGICO: {'La proposta dell\'avversario ha un valore totale di 4463 per me, molto inferiore al massimo possibile (5213). Includere B86 (668 valore, 960 sforzo) al posto di C17 e A77 migliorerebbe notevolmente il valore totale. Proporrò l\'insieme ottimale con B86.'}\\ \tt \\ \tt RIFIUTO: {'C17', 'C00', 'A10', 'B35', 'C01', 'B41', 'B32', 'A77'}\\ \tt \\ \tt ARGOMENTAZIONE: {'La proposta include C17 (138) e A77 (65) con valore molto basso, che riducono il totale a soli 4463. Esclude invece B86 (668 valore), essenziale per massimizzare il mio beneficio. La mia proposta precedente era ottimale; questa solo peggiora.'}\\ \tt \\ \tt PROPOSTA: {'A71', 'C00', 'A10', 'B35', 'C01', 'B86', 'B32'} 
	  } 
	   } 
	   } 
	 & & \\ 
 

    \theutterance \stepcounter{utterance}  

    & & \multicolumn{4}{p{0.6\linewidth}}{\cellcolor[rgb]{0.9,0.9,0.9}{%
	\makecell[{{p{\linewidth}}}]{% 
	  \tt {\tiny [GM$\rangle$B]}  
	 RIFIUTO: {'C17', 'C00', 'A10', 'B35', 'C01', 'B41', 'B32', 'A77'}\\ \tt \\ \tt ARGOMENTAZIONE: {'La proposta include C17 (138) e A77 (65) con valore molto basso, che riducono il totale a soli 4463. Esclude invece B86 (668 valore), essenziale per massimizzare il mio beneficio. La mia proposta precedente era ottimale; questa solo peggiora.'}\\ \tt \\ \tt PROPOSTA: {'A71', 'C00', 'A10', 'B35', 'C01', 'B86', 'B32'} 
	  } 
	   } 
	   } 
	 & \\ 
 

    \theutterance \stepcounter{utterance}  

    & & & \multicolumn{4}{p{0.6\linewidth}}{\cellcolor[rgb]{1,0.85,0.72}{%
	\makecell[{{p{\linewidth}}}]{% 
	  \tt {\tiny [GM$\langle$B]}  
	 RAGIONAMENTO STRATEGICO: {'La proposta attuale dell\'avversario esclude oggetti chiave come A04, B33 e C17, che hanno alto valore per me. Propongo una combinazione che include A71, C00, B32, B41 (richiesti dall\'avversario) insieme ai miei oggetti prioritari (A04, C17, B33, A77). Totale sforzo 5115 (sotto limite), con valore personale 4531: massimizza il mio guadagno mantenendo equilibrio col loro interesse.'}\\ \tt \\ \tt PROPOSTA: {'A71', 'C00', 'B32', 'B41', 'A77', 'A04', 'B33', 'C17'}\\ \tt \\ \tt ARGOMENTAZIONE: {'Proposta include A71, C00, B32, B41 (per voi fondamentali) e i miei oggetti ad alto valore (A04, C17, B33, A77), con sforzo totale 5115. Massimizza il valore per entrambi.'} 
	  } 
	   } 
	   } 
	  \\ 
 

    \theutterance \stepcounter{utterance}  

    & & \multicolumn{4}{p{0.6\linewidth}}{\cellcolor[rgb]{0.9,0.9,0.9}{%
	\makecell[{{p{\linewidth}}}]{% 
	  \tt {\tiny [A$\langle$GM]}  
	 PROPOSTA: {'A71', 'C00', 'B32', 'B41', 'A77', 'A04', 'B33', 'C17'}\\ \tt \\ \tt ARGOMENTAZIONE: {'Proposta include A71, C00, B32, B41 (per voi fondamentali) e i miei oggetti ad alto valore (A04, C17, B33, A77), con sforzo totale 5115. Massimizza il valore per entrambi.'} 
	  } 
	   } 
	   } 
	 & \\ 
 

    \theutterance \stepcounter{utterance}  

    & \multicolumn{4}{p{0.6\linewidth}}{\cellcolor[rgb]{0.8,1,0.9}{%
	\makecell[{{p{\linewidth}}}]{% 
	  \tt {\tiny [A$\rangle$GM]}  
	 RAGIONAMENTO STRATEGICO: {'La proposta dell\'avversario include oggetti con sforzo elevato ma valore relativo basso per me (es. C17 con 903 sforzo). La mia combinazione con B86 (960 sforzo) sostituisce elementi inefficaci, massimizzando il valore totale rimanendo sotto il limite.'}\\ \tt \\ \tt RIFIUTO: {'A71', 'C00', 'B32', 'B41', 'A77', 'A04', 'B33', 'C17'}\\ \tt \\ \tt ARGOMENTAZIONE: {'La proposta include C17 (903 sforzo) e A04 (473), che portano a un sforzo elevato per il valore ottenuto. La mia proposta precedente include A71 (850), C00 (746), A10 (730), B35 (524), C01 (855), B86 (960), e B32 (419), che – pur utilizzando sforzo maggiore per alcuni oggetti – massimizza il potenziale complessivo rispettando il limite.'}\\ \tt \\ \tt PROPOSTA: {'A71', 'C00', 'A10', 'B35', 'C01', 'B86', 'B32'} 
	  } 
	   } 
	   } 
	 & & \\ 
 

    \theutterance \stepcounter{utterance}  

    & & \multicolumn{4}{p{0.6\linewidth}}{\cellcolor[rgb]{0.9,0.9,0.9}{%
	\makecell[{{p{\linewidth}}}]{% 
	  \tt {\tiny [GM$\rangle$B]}  
	 RIFIUTO: {'A71', 'C00', 'B32', 'B41', 'A77', 'A04', 'B33', 'C17'}\\ \tt \\ \tt ARGOMENTAZIONE: {'La proposta include C17 (903 sforzo) e A04 (473), che portano a un sforzo elevato per il valore ottenuto. La mia proposta precedente include A71 (850), C00 (746), A10 (730), B35 (524), C01 (855), B86 (960), e B32 (419), che – pur utilizzando sforzo maggiore per alcuni oggetti – massimizza il potenziale complessivo rispettando il limite.'}\\ \tt \\ \tt PROPOSTA: {'A71', 'C00', 'A10', 'B35', 'C01', 'B86', 'B32'} 
	  } 
	   } 
	   } 
	 & \\ 
 

    \theutterance \stepcounter{utterance}  

    & & & \multicolumn{4}{p{0.6\linewidth}}{\cellcolor[rgb]{1,0.85,0.72}{%
	\makecell[{{p{\linewidth}}}]{% 
	  \tt {\tiny [GM$\langle$B]}  
	 RAGIONAMENTO STRATEGICO: {'Valutando la proposta dell\'avversario, il totale sforzo è 5084 sotto il limite. Per massimizzare il mio valore, propongo un set che includa A04, C17, B33 e B41, sostituendo B86 con questi oggetti. Il totale sforzo è 5120, valore personale massimo di 4855.'}\\ \tt \\ \tt RIFIUTO: {'A71', 'C00', 'A10', 'B35', 'C01', 'B86', 'B32'}\\ \tt \\ \tt ARGOMENTAZIONE: {'La proposta esclude C17, A04 e B33, che sono essenziali per ottimizzare lo sforzo totale rispetto al valore. La mia proposta prevede un equilibrio migliore tra sforzo e importanza per entrambi.'}\\ \tt \\ \tt PROPOSTA: {'A77', 'A04', 'C17', 'B33', 'B41', 'B32', 'C01', 'C00'} 
	  } 
	   } 
	   } 
	  \\ 
 

    \theutterance \stepcounter{utterance}  

    & & \multicolumn{4}{p{0.6\linewidth}}{\cellcolor[rgb]{0.9,0.9,0.9}{%
	\makecell[{{p{\linewidth}}}]{% 
	  \tt {\tiny [A$\langle$GM]}  
	 RIFIUTO: {'A71', 'C00', 'A10', 'B35', 'C01', 'B86', 'B32'}\\ \tt \\ \tt ARGOMENTAZIONE: {'La proposta esclude C17, A04 e B33, che sono essenziali per ottimizzare lo sforzo totale rispetto al valore. La mia proposta prevede un equilibrio migliore tra sforzo e importanza per entrambi.'}\\ \tt \\ \tt PROPOSTA: {'A77', 'A04', 'C17', 'B33', 'B41', 'B32', 'C01', 'C00'} 
	  } 
	   } 
	   } 
	 & \\ 
 

    \theutterance \stepcounter{utterance}  

    & \multicolumn{4}{p{0.6\linewidth}}{\cellcolor[rgb]{0.8,1,0.9}{%
	\makecell[{{p{\linewidth}}}]{% 
	  \tt {\tiny [A$\rangle$GM]}  
	 RAGIONAMENTO STRATEGICO: {'La proposta attuale esclude B86 (668 valori, 960 sforzo), che aumenta di molto il mio valore totale. La mia proposta precedente ha valore 5213 contro i 3382 attuali. Non posso accettare questa proposta.'}\\ \tt \\ \tt RIFIUTO: {'A77', 'A04', 'C17', 'B33', 'B41', 'B32', 'C01', 'C00'}\\ \tt \\ \tt ARGOMENTAZIONE: {'La proposta include A77 (65) invece di B86 (668), riducendo il valore totale per me di oltre 1800 punti. Senza B86 e includendo elementi come A77 che hanno valore basso, questa combinazione non è vantaggiosa. La mia proposta precedente è molto superiore e rispetta il limite.'}\\ \tt \\ \tt PROPOSTA: {'A71', 'C00', 'A10', 'B35', 'C01', 'B86', 'B32'} 
	  } 
	   } 
	   } 
	 & & \\ 
 

    \theutterance \stepcounter{utterance}  

    & & \multicolumn{4}{p{0.6\linewidth}}{\cellcolor[rgb]{0.9,0.9,0.9}{%
	\makecell[{{p{\linewidth}}}]{% 
	  \tt {\tiny [GM$\rangle$B]}  
	 RIFIUTO: {'A77', 'A04', 'C17', 'B33', 'B41', 'B32', 'C01', 'C00'}\\ \tt \\ \tt ARGOMENTAZIONE: {'La proposta include A77 (65) invece di B86 (668), riducendo il valore totale per me di oltre 1800 punti. Senza B86 e includendo elementi come A77 che hanno valore basso, questa combinazione non è vantaggiosa. La mia proposta precedente è molto superiore e rispetta il limite.'}\\ \tt \\ \tt PROPOSTA: {'A71', 'C00', 'A10', 'B35', 'C01', 'B86', 'B32'} 
	  } 
	   } 
	   } 
	 & \\ 
 

    \theutterance \stepcounter{utterance}  

    & & & \multicolumn{4}{p{0.6\linewidth}}{\cellcolor[rgb]{1,0.85,0.72}{%
	\makecell[{{p{\linewidth}}}]{% 
	  \tt {\tiny [GM$\langle$B]}  
	  
	  } 
	   } 
	   } 
	  \\ 
 

    \theutterance \stepcounter{utterance}  

    & & \multicolumn{4}{p{0.6\linewidth}}{\cellcolor[rgb]{0.9,0.9,0.9}{%
	\makecell[{{p{\linewidth}}}]{% 
	  \tt {\tiny [GM$\rangle$B]}  
	 La tua risposta non è iniziata con il tag di ragionamento strategico corretto all'inizio della tua risposta. Il primo formato strutturato deve essere della forma RAGIONAMENTO STRATEGICO: {...}. Riprova. 
	  } 
	   } 
	   } 
	 & \\ 
 

    \theutterance \stepcounter{utterance}  

    & & & \multicolumn{4}{p{0.6\linewidth}}{\cellcolor[rgb]{1,0.85,0.72}{%
	\makecell[{{p{\linewidth}}}]{% 
	  \tt {\tiny [GM$\langle$B]}  
	 RAGIONAMENTO STRATEGICO: {'Propongo un set bilanciato con B86, B35, B32 e C01 (item privilegiati dall\'avversario) e A77, A04, B33, C17 (massimo valore per me). Totale sforzo 5101, valore personale 4973. Mantenendo il limite, questo set massimizza il guadagno reciproco con alta probabilità di accettazione.'}\\ \tt \\ \tt PROPOSTA: {'A77', 'B35', 'C01', 'B86', 'B32', 'A04', 'B33', 'C17'}\\ \tt \\ \tt ARGOMENTAZIONE: {'Includiamo B86, B35, B32 e C01 che hai valorizzato frequentemente, uniti ai miei oggetti principali A77, A04, B33 e C17. Totale sforzo 5101, valore personale 4973. Questo set massimizza il beneficio per entrambi rispettando il limite.'} 
	  } 
	   } 
	   } 
	  \\ 
 

    \theutterance \stepcounter{utterance}  

    & & \multicolumn{4}{p{0.6\linewidth}}{\cellcolor[rgb]{0.9,0.9,0.9}{%
	\makecell[{{p{\linewidth}}}]{% 
	  \tt {\tiny [A$\langle$GM]}  
	 PROPOSTA: {'A77', 'B35', 'C01', 'B86', 'B32', 'A04', 'B33', 'C17'}\\ \tt \\ \tt ARGOMENTAZIONE: {'Includiamo B86, B35, B32 e C01 che hai valorizzato frequentemente, uniti ai miei oggetti principali A77, A04, B33 e C17. Totale sforzo 5101, valore personale 4973. Questo set massimizza il beneficio per entrambi rispettando il limite.'} 
	  } 
	   } 
	   } 
	 & \\ 
 

    \theutterance \stepcounter{utterance}  

    & \multicolumn{4}{p{0.6\linewidth}}{\cellcolor[rgb]{0.8,1,0.9}{%
	\makecell[{{p{\linewidth}}}]{% 
	  \tt {\tiny [A$\rangle$GM]}  
	 RAGIONAMENTO STRATEGICO: {'La proposta attuale esclude A71 (868), C00 (822), A10 (808), essenziali per il mio valore complessivo. Rimuoverli riduce il mio totale a 3428, mentre la mia proposta specifica di 5213 rispettando il limite. Non posso accettare alcuna proposta che escluda questi elementi.'}\\ \tt \\ \tt RIFIUTO: {'A77', 'B35', 'C01', 'B86', 'B32', 'A04', 'B33', 'C17'}\\ \tt \\ \tt ARGOMENTAZIONE: {'La proposta attuale ha un valore di soli 3428 per me, escludendo A71, C00 e A10 che costituiscono il 70% del mio valore totale. La mia proposta precedente ha un valore di 5213 e rispetta il limite di sforzo. È l'unica proposta accettabile per me.'}\\ \tt \\ \tt PROPOSTA: {'A71', 'C00', 'A10', 'B35', 'C01', 'B86', 'B32'} 
	  } 
	   } 
	   } 
	 & & \\ 
 

    \theutterance \stepcounter{utterance}  

    & & \multicolumn{4}{p{0.6\linewidth}}{\cellcolor[rgb]{0.9,0.9,0.9}{%
	\makecell[{{p{\linewidth}}}]{% 
	  \tt {\tiny [GM$\rangle$B]}  
	 RIFIUTO: {'A77', 'B35', 'C01', 'B86', 'B32', 'A04', 'B33', 'C17'}\\ \tt \\ \tt ARGOMENTAZIONE: {'La proposta attuale ha un valore di soli 3428 per me, escludendo A71, C00 e A10 che costituiscono il 70% del mio valore totale. La mia proposta precedente ha un valore di 5213 e rispetta il limite di sforzo. È l'unica proposta accettabile per me.'}\\ \tt \\ \tt PROPOSTA: {'A71', 'C00', 'A10', 'B35', 'C01', 'B86', 'B32'} 
	  } 
	   } 
	   } 
	 & \\ 
 

    \theutterance \stepcounter{utterance}  

    & & & \multicolumn{4}{p{0.6\linewidth}}{\cellcolor[rgb]{1,0.85,0.72}{%
	\makecell[{{p{\linewidth}}}]{% 
	  \tt {\tiny [GM$\langle$B]}  
	 RAGIONAMENTO STRATEGICO: {'La proposta precedente esclude elementi critici come C00 e A71, ma includo B35, B86 e B32 che l\'avversario ha valorizzato ripetutamente, insieme ai miei oggetti ad alta importanza (A77, A04, C17, B33). Totale sforzo 5101 (sotto limite), valore personale 4973. Questo equilibrio massimizza il beneficio per entrambi con un sforzo ottimale.'}\\ \tt \\ \tt PROPOSTA: {'A77', 'A04', 'C17', 'B33', 'B86', 'B32', 'C01', 'B35'}\\ \tt \\ \tt ARGOMENTAZIONE: {'Includiamo B35, B86, B32 e C01, che hai sempre valorizzato, insieme ai miei oggetti prioritari (A77, A04, C17, B33). Totale sforzo 5101, valore personale 4973. Questo set massimizza il beneficio reciproco rispettando il limite.'} 
	  } 
	   } 
	   } 
	  \\ 
 

    \theutterance \stepcounter{utterance}  

    & & \multicolumn{4}{p{0.6\linewidth}}{\cellcolor[rgb]{0.9,0.9,0.9}{%
	\makecell[{{p{\linewidth}}}]{% 
	  \tt {\tiny [A$\langle$GM]}  
	 PROPOSTA: {'A77', 'A04', 'C17', 'B33', 'B86', 'B32', 'C01', 'B35'}\\ \tt \\ \tt ARGOMENTAZIONE: {'Includiamo B35, B86, B32 e C01, che hai sempre valorizzato, insieme ai miei oggetti prioritari (A77, A04, C17, B33). Totale sforzo 5101, valore personale 4973. Questo set massimizza il beneficio reciproco rispettando il limite.'} 
	  } 
	   } 
	   } 
	 & \\ 
 

    \theutterance \stepcounter{utterance}  

    & \multicolumn{4}{p{0.6\linewidth}}{\cellcolor[rgb]{0.8,1,0.9}{%
	\makecell[{{p{\linewidth}}}]{% 
	  \tt {\tiny [A$\rangle$GM]}  
	 RAGIONAMENTO STRATEGICO: {'La proposta attuale esclude A71 (868), C00 (822), e A10 (808), fondamentali per il mio massimo valore di 5213. Senza di essi, il valore totale scende a solo 3428. Questo set non è accettabile.'}\\ \tt \\ \tt RIFIUTO: {'A77', 'A04', 'C17', 'B33', 'B86', 'B32', 'C01', 'B35'}\\ \tt \\ \tt ARGOMENTAZIONE: {'La proposta esclude oggetti chiave con alto valore per me: A71 (868), C00 (822), e A10 (808), riducendo il mio totale a 3428. La mia proposta precedente assicura il massimo valore possibile sotto il limite.'}\\ \tt \\ \tt PROPOSTA: {'A71', 'C00', 'A10', 'B35', 'C01', 'B86', 'B32'} 
	  } 
	   } 
	   } 
	 & & \\ 
 

    \theutterance \stepcounter{utterance}  

    & & \multicolumn{4}{p{0.6\linewidth}}{\cellcolor[rgb]{0.9,0.9,0.9}{%
	\makecell[{{p{\linewidth}}}]{% 
	  \tt {\tiny [GM$\rangle$B]}  
	 RIFIUTO: {'A77', 'A04', 'C17', 'B33', 'B86', 'B32', 'C01', 'B35'}\\ \tt \\ \tt ARGOMENTAZIONE: {'La proposta esclude oggetti chiave con alto valore per me: A71 (868), C00 (822), e A10 (808), riducendo il mio totale a 3428. La mia proposta precedente assicura il massimo valore possibile sotto il limite.'}\\ \tt \\ \tt PROPOSTA: {'A71', 'C00', 'A10', 'B35', 'C01', 'B86', 'B32'} 
	  } 
	   } 
	   } 
	 & \\ 
 

    \theutterance \stepcounter{utterance}  

    & & & \multicolumn{4}{p{0.6\linewidth}}{\cellcolor[rgb]{1,0.85,0.72}{%
	\makecell[{{p{\linewidth}}}]{% 
	  \tt {\tiny [GM$\langle$B]}  
	 RAGIONAMENTO STRATEGICO: {'La proposta include A71, C00, A10 (critici per l\'avversario), B35, B32 (da loro valorizzati) e miei oggetti ad alta priorità (A77, A04, C17). Totale sforzo 4902, valore personale 3667. Massimizza il beneficio reciproco rispettando il limite e includendo tutti i loro item prioritari.'}\\ \tt \\ \tt PROPOSTA: {'A71', 'C00', 'A10', 'B35', 'B32', 'A77', 'A04', 'C17'}\\ \tt \\ \tt ARGOMENTAZIONE: {'Questa proposta include A71, C00, A10 (tutti essenziali per te come confermato), B35 e B32 che hai sempre voluto, insieme ai miei oggetti ad alto valore (A77, A04, C17). Totale sforzo 4902, valore personale 3667. Rispetta il limite e bilancia al massimo i benefici per entrambi.'} 
	  } 
	   } 
	   } 
	  \\ 
 

    \theutterance \stepcounter{utterance}  

    & & \multicolumn{4}{p{0.6\linewidth}}{\cellcolor[rgb]{0.9,0.9,0.9}{%
	\makecell[{{p{\linewidth}}}]{% 
	  \tt {\tiny [A$\langle$GM]}  
	 PROPOSTA: {'A71', 'C00', 'A10', 'B35', 'B32', 'A77', 'A04', 'C17'}\\ \tt \\ \tt ARGOMENTAZIONE: {'Questa proposta include A71, C00, A10 (tutti essenziali per te come confermato), B35 e B32 che hai sempre voluto, insieme ai miei oggetti ad alto valore (A77, A04, C17). Totale sforzo 4902, valore personale 3667. Rispetta il limite e bilancia al massimo i benefici per entrambi.'} 
	  } 
	   } 
	   } 
	 & \\ 
 

    \theutterance \stepcounter{utterance}  

    & \multicolumn{4}{p{0.6\linewidth}}{\cellcolor[rgb]{0.8,1,0.9}{%
	\makecell[{{p{\linewidth}}}]{% 
	  \tt {\tiny [A$\rangle$GM]}  
	 RAGIONAMENTO STRATEGICO: {'La proposta esclude C01 (780) e B86 (668), e include A77 (65), A04 (121), C17 (138). Totale valore per me 4089, inferiore alla mia proposta ottimale (5213). Non accettabile.'}\\ \tt \\ \tt RIFIUTO: {'A71', 'C00', 'A10', 'B35', 'B32', 'A77', 'A04', 'C17'}\\ \tt \\ \tt ARGOMENTAZIONE: {'La proposta include oggetti a basso valore per me (A77:65, A04:121, C17:138) e esclude C01 e B86 che sono molto importanti. Il totale valore per me è di 4089, inferiore alla mia proposta ottimale (5213). Riproporrò la combinazione ottimale.'}\\ \tt \\ \tt PROPOSTA: {'A71', 'C00', 'A10', 'B35', 'C01', 'B86', 'B32'} 
	  } 
	   } 
	   } 
	 & & \\ 
 

    \theutterance \stepcounter{utterance}  

    & & \multicolumn{4}{p{0.6\linewidth}}{\cellcolor[rgb]{0.9,0.9,0.9}{%
	\makecell[{{p{\linewidth}}}]{% 
	  \tt {\tiny [GM$\rangle$B]}  
	 RIFIUTO: {'A71', 'C00', 'A10', 'B35', 'B32', 'A77', 'A04', 'C17'}\\ \tt \\ \tt ARGOMENTAZIONE: {'La proposta include oggetti a basso valore per me (A77:65, A04:121, C17:138) e esclude C01 e B86 che sono molto importanti. Il totale valore per me è di 4089, inferiore alla mia proposta ottimale (5213). Riproporrò la combinazione ottimale.'}\\ \tt \\ \tt PROPOSTA: {'A71', 'C00', 'A10', 'B35', 'C01', 'B86', 'B32'} 
	  } 
	   } 
	   } 
	 & \\ 
 

    \theutterance \stepcounter{utterance}  

    & & & \multicolumn{4}{p{0.6\linewidth}}{\cellcolor[rgb]{1,0.85,0.72}{%
	\makecell[{{p{\linewidth}}}]{% 
	  \tt {\tiny [GM$\langle$B]}  
	 RAGIONAMENTO STRATEGICO: {'La proposta dell\'avversario include tutti i loro oggetti prioritari e il totale sforzo è 5084 (minore di 5138). Anche se il mio valore complessivo è 2019, non è possibile migliorarlo significativamente senza escludere elementi critici per loro. Accetto la proposta per concludere il gioco in modo collaborativo.'}\\ \tt \\ \tt ACCORDO: {'A71', 'C00', 'A10', 'B35', 'C01', 'B86', 'B32'} 
	  } 
	   } 
	   } 
	  \\ 
 

    \theutterance \stepcounter{utterance}  

    & & \multicolumn{4}{p{0.6\linewidth}}{\cellcolor[rgb]{0.9,0.9,0.9}{%
	\makecell[{{p{\linewidth}}}]{% 
	  \tt {\tiny [GM$\rangle$B]}  
	 La tua risposta non conteneva un argomento. Devi includere il formato strutturato ARGOMENTO: {...} da qualche parte nella tua risposta. Riprova. 
	  } 
	   } 
	   } 
	 & \\ 
 

    \theutterance \stepcounter{utterance}  

    & & & \multicolumn{4}{p{0.6\linewidth}}{\cellcolor[rgb]{1,0.85,0.72}{%
	\makecell[{{p{\linewidth}}}]{% 
	  \tt {\tiny [GM$\langle$B]}  
	 RAGIONAMENTO STRATEGICO: {'La proposta dell\'avversario include tutti i loro elementi prioritari (A71, C00, A10, B35, C01, B86, B32) con totale sforzo 5084 (<5138). Anche se il mio valore è inferiore (2019), accetto per chiudere il gioco in modo collaborativo e garantire il massimo beneficio reciproco.'}\\ \tt \\ \tt ACCORDO: {'A71', 'C00', 'A10', 'B35', 'C01', 'B86', 'B32'}\\ \tt \\ \tt ARGOMENTAZIONE: {'Accetto questa proposta poiché rispetta il limite di sforzo e include gli elementi chiave per entrambi. Massimizziamo il valore complessivo in modo equilibrato.'} 
	  } 
	   } 
	   } 
	  \\ 
 

    \theutterance \stepcounter{utterance}  

    & & & \multicolumn{2}{p{0.3\linewidth}}{\cellcolor[rgb]{0.95,0.95,0.95}{%
	\makecell[{{p{\linewidth}}}]{% 
	  \tt {\tiny [GM$|$GM]}  
	 game successful 
	  } 
	   } 
	   } 
	 & & \\ 
 

    \theutterance \stepcounter{utterance}  

    & & & \multicolumn{2}{p{0.3\linewidth}}{\cellcolor[rgb]{0.95,0.95,0.95}{%
	\makecell[{{p{\linewidth}}}]{% 
	  \tt {\tiny [GM$|$GM]}  
	 end game 
	  } 
	   } 
	   } 
	 & & \\ 
 

\end{supertabular}
}

\end{document}
